\documentclass[11pt]{article}       % set main text size
\usepackage[letterpaper,            % set paper size to letterpaper. change to a4paper for resumes outside of North America
top=0.5in,                          % specify top page margin
bottom=0.5in,                       % specify bottom page margin
left=0.5in,                         % specify left page margin
right=0.5in]{geometry}              % specify right page margin
                       
\usepackage{XCharter}               % set font
\usepackage[T1]{fontenc}            % output encoding
\usepackage[utf8]{inputenc}         % input encoding
\usepackage{enumitem}               % enable lists for bullet points: itemize and \item
\usepackage[colorlinks=true, linkcolor=black, urlcolor=blue]{hyperref}    % format hyperlinks
\usepackage{titlesec}               % enable section title customization
\raggedright                        % disable text justification
\pagestyle{empty}                   % disable page numbering
\usepackage{comment}                % block comment

% ensure PDF output will be all-Unicode and machine-readable
\input{glyphtounicode}
\pdfgentounicode=1

% format section headings: bolding, size, white space above and below
\titleformat{\section}{\bfseries\large}{}{0pt}{}[\vspace{1pt}\titlerule\vspace{-6.5pt}]

% format bullet points: size, white space above and below, white space between bullets
\renewcommand\labelitemi{$\vcenter{\hbox{\small$\bullet$}}$}
\setlist[itemize]{itemsep=-2pt, leftmargin=12pt, topsep=6pt} %%% Test various topsep values to fix vertical spacing errors

% resume starts here
\begin{document}

% name
\centerline{\Huge Thuy Nguyen}

\vspace{5pt}

% contact information
\centerline{\href{mailto:thuybohr@gmail.com}{thuybohr@gmail.com} | 206-423-9827 |
\href{https://www.linkedin.com/in/thuy-nguyen-46505121b/}{linkedin.com/in/thuy-nguyen-46505121b}}

\centerline{\href{https://github.com/Monica20030707}{github.com/Monica20030707} | 
\href{https://monica20030707.github.io/}{monica20030707.github.io/}} 

\vspace{-9pt}

% education section
\section*{Education}
\textbf{Bellevue College} -- BS in Computer Science -- 3.6 GPA \hfill September 2021 -- June 2025\\
\vspace{-9pt}
\begin{itemize}
  \item Award: Amazon Scholarship, Activities: Social Media Manager of Computer Science Club
\end{itemize}

\vspace{-18.5pt}

% experience section
\section*{Experience}

\textbf{Associate Full Stack Developer,} {Forge Tek Studios} -- Remote \hfill August 2025 -- Present 
\vspace{-9pt}
\begin{itemize}
  \item Architected a modular React frontend with a serverless backend, leveraging custom hooks and context for efficient state management and performance-sensitive UI updates.
  \item Optimized client-side rendering and data synchronization, reducing load and ingestion overhead by 80\% for production-ready web and mobile applications.
  \item Implemented dynamic data fetching and caching strategies for third-party APIs like Steam and Fab Marketplace to reduce latency and improve user experience.
  \item Owned the end-to-end development and deployment lifecycle, ensuring high standards for code quality and performance using serverless technologies.
\end{itemize}

\textbf{Full Stack Developer Intern,} {METY Technology} -- Remote \hfill September 2024 -- June 2025 
\vspace{-9pt}
\begin{itemize}
  \item Integrated real-time machine learning models (MoveNet) with React to enable low-latency, gesture-based user interaction directly in the browser.
  \item Developed secure and scalable RESTful APIs with Express.js and Node.js with a robust state management layer using Redux Toolkit.
  \item Automated testing and deployment processes via a custom CI/CD pipeline on AWS Amplify, ensuring consistent delivery and quality standards.
  \item Enhanced user engagement by 25\% through the development of highly responsive and intuitive real-time visual feedback and state managing.
\end{itemize}

\textbf{Open-Source Intern,} {CodeDay} -- Seattle, WA \hfill April 2025 -- May 2025 \\
\vspace{-9pt}
\begin{itemize}
  \item Orchestrated Docker-based development environments to ensure consistent application behavior across local and staging stages.
  \item Authored comprehensive database integration tests to resolve high-priority open-source tickets and prevent regression in an Agile environment.
  \item Drove project velocity as Scrum Master, meeting all client deadlines through proactive blocker removal and effective team communication.
  \item Facilitated daily stand-ups and retrospectives, aligning a cross-functional team on weekly sprint goals and technical documentation.
\end{itemize}

\textbf{Coding Competition Developer,} {Bellevue College} -- Bellevue, WA \hfill April 2024 -- May 2025 \\
\vspace{-9pt}
\begin{itemize}
  \item Built a dynamic leaderboard system using PHP and SQL for real-time performance tracking during competitive coding events.
  \item Authored automated test cases across 5 languages (including Java and Python) to validate participant submissions for algorithm accuracy and efficiency.
  \item Enhanced web interfaces with HTML and CSS, boosting user satisfaction by 30\% through improved UI/UX and interface responsiveness.
  \item Directed onsite technical operations for 30+ participants, ensuring a seamless competition experience and technical stability.
\end{itemize}

\vspace{-18.5pt}

% projects section
\section*{Projects}

\textbf{Stock Ratio Tracker}{ - Python, React, Typescript} \hfill \href{https://github.com/Monica20030707/tradingDashboard-UI}{github.com} \\
\vspace{-9pt}
\begin{itemize}
  \item Architected the data-flow pipeline connecting Python backend calculations to a live React dashboard, ensuring stable, low-latency ratio updates. 
  \item Used JPMorgan Chase’s open-source library called Perspective to generate a live graph that displays a data feed in a clear and visually appealing way.
  \item Implemented automated unit tests to ensure correctness and reliability of real-time financial data visualization.
\end{itemize}

\textbf{Reverse-Polish Calculator}{ - Java, ANTLR, JUnit, OOP Concepts} \hfill \href{https://github.com/Monica20030707/reverse-Polish_calculator}{github.com} \\
\vspace{-9pt}
\begin{itemize}
  \item Designed and implemented a grammar-driven parser using ANTLR and the Visitor design pattern for abstract syntax tree traversal, reinforcing OOP principles.
  \item Leveraged stacks and tree-based data structures to correctly parse and compute complex nested reverse-polish notation expressions.
  \item Achieved 100\% test coverage for all arithmetic edge cases using JUnit-driven automated testing to ensure system reliability.
\end{itemize}

\vspace{-18.5pt}

% skills section
\section*{Skills}
\textbf{Languages:} Java, Python, JavaScript, TypeScript, SQL (MySQL), PHP, HTML, CSS, Kotlin, Go, C++\\
\textbf{Skills:} Data Structures, Object-Oriented Programming (OOP), Unit Testing (JUnit, Jest), Agile/Scrum, RESTful APIs, Networking Basics, CI/CD, Code Reviews\\
\textbf{Tools:} React, React Native, Android, iOS, Redux Toolkit, Node.js, Express.js, Docker, AWS, Git, ANTLR, JUnit\\

\vspace{-6.5pt}

\end{document}
